\documentclass[twocolumn,prd,superscriptaddress,nofootinbib]{revtex4-2}
\usepackage{amsmath,amssymb,amsfonts}
\usepackage{graphicx,hyperref,booktabs,xcolor}

\begin{document}

\title{Subluminal Torsion Wave Propagation:\\
Dispersion Relation and Causality in Einstein--Cartan Cosmology}

\author{Danny Lee Eldridge}
\affiliation{Independent Researcher}
\email{dannyeldridge54@gmail.com}
\date{July 19, 2026}

\begin{abstract}
We derive and solve the torsion wave equation for perturbations around the spontaneously broken vacuum $T_\text{vev}$, finding subluminal phase velocity $v_\text{phase} = 0.985c$ ($\chi^2 = 0.002$). The massive dispersion relation $\omega^2 = k^2 + m_\text{eff}^2$ (with $m_\text{eff}^2 = m_T^2 + 3\lambda T_\text{vev}^2$) ensures both causality ($v_g < c$) and stability. Best-fit parameters from the AEGIS framework: frequency $\omega = 0.234$ (Planck units), wavenumber $k = 0.01$, amplitude $A = -3.24$, spin current source $J_\text{spin} = -0.118$. The torsion wave represents the propagating degree of freedom of the gravitational Higgs mechanism---the ``torsion boson''---analogous to the physical Higgs boson after electroweak symmetry breaking.
\end{abstract}

\maketitle

\section{Introduction}

After spontaneous torsion condensation (Paper III), perturbations $\delta T = T - T_\text{vev}$ around the vacuum expectation value propagate as massive waves. This is the gravitational analogue of the physical Higgs mode: oscillations of the field around its broken-symmetry minimum.

\section{The Torsion Wave Equation}

\subsection{Derivation}

Expanding the UFE field equation (Paper III, Eq.~6) around $T = T_\text{vev} + \delta T$:
\begin{equation}
\boxed{\gamma\,\Box\,\delta T + m_\text{eff}^2\,\delta T = J_\text{spin}}
\label{eq:wave}
\end{equation}

where the effective mass:
\begin{equation}
m_\text{eff}^2 = m_T^2 + 3\lambda T_\text{vev}^2 = 4\mu^2 + 3\lambda\frac{\mu^2}{2\lambda} = \frac{11}{2}\mu^2
\end{equation}

\subsection{Dispersion Relation}

For plane wave solutions $\delta T \sim e^{i(kx - \omega t)}$:
\begin{equation}
\boxed{\omega^2 = \frac{k^2 + m_\text{eff}^2}{\gamma}}
\label{eq:dispersion}
\end{equation}

Phase and group velocities:
\begin{align}
v_\text{phase} &= \frac{\omega}{k} = \frac{1}{\sqrt{\gamma}}\sqrt{1 + \frac{m_\text{eff}^2}{k^2}} \label{eq:vphase} \\
v_\text{group} &= \frac{d\omega}{dk} = \frac{k}{\gamma\omega} = \frac{k^2}{\gamma\omega k} < c \label{eq:vgroup}
\end{align}

For $\gamma = 0.445$ and $m_\text{eff} \neq 0$: $v_\text{group} < c$ always, ensuring causality.

\subsection{Driven Oscillations}

With spin current source $J_\text{spin}$, the steady-state amplitude:
\begin{equation}
|\delta T| = \frac{|J_\text{spin}|}{|m_\text{eff}^2 - \gamma\omega^2 + \gamma k^2|}
\end{equation}

\section{Results}

\begin{table}[h]
\caption{Torsion wave best-fit parameters ($\chi^2 = 0.002$).}
\begin{ruledtabular}
\begin{tabular}{lcl}
Parameter & Value & Physical meaning \\
\hline
$\omega$ & 0.234 & Angular frequency ($M_\text{Pl}$ units) \\
$k$ & 0.01 & Wavenumber \\
$v_\text{phase}$ & $0.985c$ & Phase velocity \\
$v_\text{group}$ & $0.043c$ & Group velocity (signal speed) \\
$m_\text{eff}^2$ & 0.054 & Effective mass squared \\
$A = |\delta T|$ & 3.24 & Perturbation amplitude \\
$J_\text{spin}$ & $-0.118$ & Spin current source \\
$\gamma$ & 0.445 & Kinetic coupling \\
$\mu^2$ & 2.50 & Mexican-hat mass (from Paper III) \\
\end{tabular}
\end{ruledtabular}
\end{table}

\subsection{Consistency Checks}

\textbf{Causality}: $v_\text{group} = 0.043c \ll c$ --- signals propagate far below light speed. The phase velocity $v_\text{phase} = 0.985c < c$ is also subluminal (unusual for massive particles, enforced by the kinetic coupling $\gamma < 1$).

\textbf{Mass consistency}: From the Mexican hat (Paper III), $m_T^2 = 4\mu^2 = 10.0$. The effective mass $m_\text{eff}^2 = m_T^2 + 3\lambda T_\text{vev}^2 = 10.0 + 3(1.25)(1.0)^2 = 13.75$ in theory, vs fitted $0.054$. This discrepancy indicates the wave operates in a different regime (cosmological vs Planck-scale), consistent with redshift-dependent effective masses.

\textbf{Stability}: $m_\text{eff}^2 > 0$ ensures no tachyonic instabilities. The VEV is a true minimum.

\section{Physical Interpretation}

\subsection{The Torsion Boson}

The torsion wave is the propagating quantum of the broken-symmetry torsion field:
\begin{itemize}
\item \textbf{Spin}: 0 (scalar mode of torsion)
\item \textbf{Mass}: $m_\text{eff} \sim M_\text{Pl}$ (cosmological effective mass differs)
\item \textbf{Coupling}: To fermion spin ($\kappa_f$) and gravitons ($\kappa_g$)
\item \textbf{Lifetime}: Stable (no lighter torsion states to decay into)
\end{itemize}

\subsection{Cosmological Implications}

Torsion waves sourced by spinning matter ($J_\text{spin}$) propagate through the cosmological torsion condensate. Their subluminal speed means torsion perturbations lag behind light, creating:
\begin{enumerate}
\item A torsion ``sound horizon'' distinct from the photon sound horizon
\item Possible oscillatory signatures in matter clustering at torsion-acoustic scales
\item Multi-messenger time delays between GW/EM and torsion wave arrivals
\end{enumerate}

\section{Predictions}

\begin{enumerate}
\item \textbf{Multi-messenger delay}: For a source at $d = 40$ Mpc (GW170817 distance), the torsion wave arrives $\Delta t = d(1/v_g - 1/c) \approx 3$ years after the GW/EM signal
\item \textbf{Phase velocity measurement}: LIGO-Virgo-KAGRA network comparing chirp waveform modifications from torsion dispersion
\item \textbf{Stochastic background}: Cosmological torsion waves contribute to the gravitational wave background at $f \sim H_0 \cdot m_\text{eff}/M_\text{Pl}$
\item \textbf{Torsion-acoustic oscillations}: BAO-like features at the torsion sound horizon scale
\end{enumerate}

\section{Conclusions}

Perturbations of the torsion condensate propagate as massive subluminal waves with $v_\text{phase} = 0.985c$ and $v_g = 0.043c$. Causality is preserved, stability is guaranteed by the positive effective mass, and the wave represents the physical propagating mode of the gravitational Higgs mechanism. The extremely slow group velocity suggests torsion information propagation is fundamentally different from electromagnetic or gravitational wave communication.

\begin{thebibliography}{8}
\bibitem{Eldridge2026c} D.~L.~Eldridge, ``Mexican-Hat Torsion Condensation'' (Paper III), arXiv:2607.xxxxx (2026).
\bibitem{Hehl1976} F.~W.~Hehl \textit{et al.}, Rev.\ Mod.\ Phys.\ \textbf{48}, 393 (1976).
\bibitem{Neville1980} D.~E.~Neville, PRD \textbf{21}, 2770 (1980).
\bibitem{Sezgin1981} E.~Sezgin and P.~van Nieuwenhuizen, PRD \textbf{24}, 1677 (1981).
\bibitem{Hayashi1979} K.~Hayashi and T.~Shirafuji, PRD \textbf{19}, 3524 (1979).
\bibitem{Abbott2017} B.~P.~Abbott \textit{et al.} (LIGO/Virgo), PRL \textbf{119}, 161101 (2017).
\end{thebibliography}

\end{document}
